\documentclass[11pt]{article}

\usepackage[margin=1in]{geometry}
\usepackage[T1]{fontenc}
\usepackage{lmodern}
\usepackage{microtype}
\usepackage{amsmath, amssymb, amsthm, mathtools}
\usepackage{enumitem}
\usepackage{booktabs}
\usepackage{longtable}
\usepackage{array}
\usepackage{xcolor}
\usepackage{hyperref}
\usepackage{fancyhdr}
\usepackage{titlesec}

\hypersetup{
    colorlinks=true,
    linkcolor=blue!60!black,
    urlcolor=blue!60!black,
    citecolor=blue!60!black,
    pdftitle={Math 617: Number Theory I},
    pdfauthor={Johns Hopkins University}
}

\pagestyle{fancy}
\fancyhf{}
\lhead{Math 617: Number Theory I}
\rhead{Fall 2023}
\cfoot{\thepage}

\setlist[itemize]{noitemsep, topsep=3pt}
\setlist[enumerate]{noitemsep, topsep=3pt}

\newtheorem{definition}{Definition}[section]
\newtheorem{theorem}[definition]{Theorem}
\newtheorem{proposition}[definition]{Proposition}
\newtheorem{lemma}[definition]{Lemma}
\newtheorem{example}[definition]{Example}
\newtheorem{remark}[definition]{Remark}

\titleformat{\section}
  {\large\bfseries}
  {\thesection.}
  {0.5em}
  {}

\titleformat{\subsection}
  {\normalsize\bfseries}
  {\thesubsection}
  {0.5em}
  {}

\begin{document}

\begin{center}
    {\Large \textbf{Math 617: Number Theory I}}\\[0.4em]
    {\large Fall Semester 2023}\\[0.4em]
    August 28, 2023--December 8, 2023
\end{center}

\vspace{1em}

\begin{center}
\begin{tabular}{ll}
\textbf{Instructor:} & Huajie Li \\
\textbf{Email:} & \href{mailto:hli213@jhu.edu}{hli213@jhu.edu} \\
\textbf{Office:} & Krieger 220 \\
\textbf{Lectures:} & Tuesdays and Thursdays, 9:00--10:15 AM \\
\textbf{Location:} & Hodson 315 \\
\textbf{Office Hours:} & Tuesdays, 10:30--11:30 AM, or by appointment \\
\end{tabular}
\end{center}

\section{Course Description}

Math 617 is a graduate course introducing the foundations of modern algebraic
and arithmetic number theory. The course develops the arithmetic of number
fields and their rings of integers, studies ramification and completions, and
introduces the global--local viewpoint through global fields, local fields,
adèles, and idèles. If time permits, the course concludes with selected topics
from Tate's thesis.

The course emphasizes proofs, structural understanding, and problem solving.
Homework is an essential component of the course and may include both routine
exercises and problems requiring substantial theoretical insight.

\section{Course Information}

\subsection{Prerequisites}

Students are expected to have graduate-level familiarity with algebra and
analysis, approximately at the level of qualifying examinations.

Important prerequisite topics include:

\begin{itemize}
    \item Groups, rings, ideals, modules, and fields.
    \item Polynomial rings and field extensions.
    \item Vector spaces, linear maps, and tensor products.
    \item Noetherian and Dedekind domains.
    \item Basic Galois theory.
    \item Integral extensions and localization.
    \item Basic homological and commutative algebra terminology.
    \item Metric spaces, normed spaces, and complete spaces.
    \item Basic point-set topology.
    \item Measure and integration at an introductory graduate level.
    \item Infinite series, products, and convergence.
    \item Basic complex analysis and Fourier analysis, especially for the
          optional material on Tate's thesis.
\end{itemize}

\subsection{References}

There is no required textbook. The following references provide standard
background and complementary treatments:

\begin{itemize}
    \item J. Neukirch, \emph{Algebraic Number Theory}.
    \item A. Weil, \emph{Basic Number Theory}.
    \item S. Lang, \emph{Algebraic Number Theory}.
    \item J. W. S. Cassels and A. Fr\"ohlich, editors,
          \emph{Algebraic Number Theory}.
\end{itemize}

Additional references and lecture notes may be introduced during the semester.

\section{Course Structure}

The course is organized into the following modules:

\begin{enumerate}
    \item Algebraic preliminaries and arithmetic of number fields.
    \item Dedekind domains and factorization of ideals.
    \item Ramification and discriminants.
    \item Valuations, completions, and local fields.
    \item Global fields and local--global principles.
    \item Adèles and idèles.
    \item Selected topics related to Tate's thesis, if time permits.
\end{enumerate}

The schedule is approximate. The instructor may adjust the amount of time
devoted to each module depending on the pace of the class and the level of
discussion.

\section{Learning Objectives}

By the end of the course, students should be able to:

\begin{enumerate}
    \item Work fluently with rings of integers in number fields.
    \item Prove and use the structure theorem for ideals in Dedekind domains.
    \item Explain the difference between factorization of elements and
          factorization of ideals.
    \item Analyze the behavior of prime ideals under extensions of number
          fields.
    \item Define and compute ramification indices and residue degrees.
    \item Relate discriminants, differents, and ramification.
    \item Use valuations to describe arithmetic and analytic properties of
          fields.
    \item Construct completions of fields with respect to absolute values.
    \item Explain the relationship between global fields and their local
          completions.
    \item Work with finite extensions of local fields.
    \item Define the ring of adèles and the group of idèles.
    \item Understand why adèles and idèles provide a useful global--local
          framework.
    \item Formulate and apply basic local--global principles.
    \item Read and construct rigorous proofs in algebraic number theory.
    \item Solve unfamiliar problems independently and communicate mathematical
          arguments clearly.
\end{enumerate}

\section{Detailed Course Modules}

\subsection{Module 1: Algebraic and Arithmetic Preliminaries}

\textbf{Suggested duration:} Weeks 1--2

\subsubsection*{Main themes}

\begin{itemize}
    \item Number fields and algebraic extensions.
    \item Rings of integers.
    \item Integral elements and integral closure.
    \item Norms and traces.
    \item Embeddings and signatures.
    \item Fractional ideals and localization.
\end{itemize}

\subsubsection*{Suggested chapters}

\begin{enumerate}
    \item Number fields and their embeddings.
    \item Integral dependence and rings of integers.
    \item Norms, traces, and discriminants.
    \item Fractional ideals and localization.
\end{enumerate}

\subsubsection*{Important concepts}

\begin{itemize}
    \item Number field \(K\).
    \item Ring of integers \(\mathcal{O}_K\).
    \item Minimal polynomial.
    \item Integral element.
    \item Field norm \(N_{K/\mathbb{Q}}(\alpha)\).
    \item Field trace \(\operatorname{Tr}_{K/\mathbb{Q}}(\alpha)\).
    \item Fractional ideal.
    \item Localization \(S^{-1}R\).
    \item Discriminant of a basis.
    \item Real and complex embeddings.
\end{itemize}

\subsubsection*{Learning objectives}

Students should be able to:

\begin{itemize}
    \item Determine whether an element is integral over a given ring.
    \item Describe the ring of integers of basic number fields.
    \item Compute norms and traces using multiplication matrices.
    \item Explain the role of embeddings in the arithmetic of a number field.
    \item Work with fractional ideals and localizations.
\end{itemize}

\subsection{Module 2: Dedekind Domains and Ideal Factorization}

\textbf{Suggested duration:} Weeks 3--5

\subsubsection*{Main themes}

\begin{itemize}
    \item Definition and characterization of Dedekind domains.
    \item Unique factorization of nonzero ideals.
    \item Prime ideals and maximal ideals.
    \item Fractional ideal groups.
    \item Principal ideals and the ideal class group.
    \item Failure of unique factorization of elements.
\end{itemize}

\subsubsection*{Suggested chapters}

\begin{enumerate}
    \item Dedekind domains.
    \item Factorization of ideals.
    \item The ideal class group.
    \item Examples from quadratic number fields.
\end{enumerate}

\subsubsection*{Important concepts}

\begin{itemize}
    \item Noetherian domain.
    \item Integrally closed domain.
    \item Krull dimension one.
    \item Dedekind domain.
    \item Nonzero prime ideal.
    \item Invertible ideal.
    \item Fractional ideal group.
    \item Principal ideal.
    \item Ideal class group \(\operatorname{Cl}(\mathcal{O}_K)\).
    \item Class number.
\end{itemize}

\subsubsection*{Core theorem}

Every nonzero proper ideal of a Dedekind domain factors uniquely as a product
of nonzero prime ideals:
\[
    \mathfrak{a}
    =
    \mathfrak{p}_1^{e_1}\cdots\mathfrak{p}_r^{e_r}.
\]

\subsubsection*{Learning objectives}

Students should be able to:

\begin{itemize}
    \item Verify that a ring is a Dedekind domain.
    \item Factor ideals into prime ideals.
    \item Prove uniqueness of ideal factorization.
    \item Distinguish element factorization from ideal factorization.
    \item Compute or estimate class groups in elementary examples.
    \item Explain how the class group measures the failure of unique
          factorization of elements.
\end{itemize}

\subsection{Module 3: Ramification and Discriminants}

\textbf{Suggested duration:} Weeks 5--7

\subsubsection*{Main themes}

\begin{itemize}
    \item Extension of prime ideals.
    \item Factorization of primes in extensions.
    \item Ramification indices and residue degrees.
    \item Decomposition groups and inertia groups.
    \item Discriminants and differents.
    \item Dedekind's factorization criterion.
\end{itemize}

\subsubsection*{Suggested chapters}

\begin{enumerate}
    \item Extension of ideals.
    \item Splitting of primes.
    \item Ramification.
    \item Discriminants and the different.
    \item Galois-theoretic interpretation of decomposition.
\end{enumerate}

\subsubsection*{Important concepts}

For a prime ideal \(\mathfrak{p}\) of \(\mathcal{O}_K\), its extension to
\(\mathcal{O}_L\) may factor as
\[
    \mathfrak{p}\mathcal{O}_L
    =
    \mathfrak{P}_1^{e_1}\cdots\mathfrak{P}_g^{e_g}.
\]

The associated quantities are:

\begin{itemize}
    \item \(e_i\): ramification index.
    \item \(f_i=[\mathcal{O}_L/\mathfrak{P}_i:
          \mathcal{O}_K/\mathfrak{p}]\): residue degree.
    \item \(g\): number of primes above \(\mathfrak{p}\).
\end{itemize}

In the separable case, these satisfy
\[
    \sum_{i=1}^{g}e_i f_i=[L:K].
\]

Other central concepts include:

\begin{itemize}
    \item Unramified, tamely ramified, and wildly ramified primes.
    \item Decomposition group \(D(\mathfrak{P}/\mathfrak{p})\).
    \item Inertia group \(I(\mathfrak{P}/\mathfrak{p})\).
    \item Different ideal.
    \item Discriminant ideal.
    \item Residue field extension.
\end{itemize}

\subsubsection*{Learning objectives}

Students should be able to:

\begin{itemize}
    \item Factor primes in basic extensions of number fields.
    \item Compute ramification indices and residue degrees.
    \item Apply the relation
          \(\sum e_i f_i=[L:K]\).
    \item Identify ramified primes using discriminants.
    \item Explain ramification through both ideal-theoretic and
          Galois-theoretic language.
    \item Distinguish decomposition, inertia, and ramification.
\end{itemize}

\subsection{Module 4: Valuations, Completions, and Local Fields}

\textbf{Suggested duration:} Weeks 7--9

\subsubsection*{Main themes}

\begin{itemize}
    \item Absolute values and valuations.
    \item Nonarchimedean absolute values.
    \item Completion of \(\mathbb{Q}\) and number fields.
    \item \(p\)-adic numbers.
    \item Local fields.
    \item Extensions of local fields.
    \item Hensel's lemma.
\end{itemize}

\subsubsection*{Suggested chapters}

\begin{enumerate}
    \item Valued fields.
    \item The \(p\)-adic numbers.
    \item Completions of number fields.
    \item Finite extensions of local fields.
    \item Hensel's lemma and lifting arguments.
\end{enumerate}

\subsubsection*{Important concepts}

\begin{itemize}
    \item Absolute value \(|\cdot|\).
    \item Valuation \(v\).
    \item Nonarchimedean valuation.
    \item Ultrametric inequality.
    \item Completion \(K_v\).
    \item \(p\)-adic field \(\mathbb{Q}_p\).
    \item Ring of integers \(\mathcal{O}_{K_v}\).
    \item Maximal ideal \(\mathfrak{m}_{K_v}\).
    \item Residue field.
    \item Uniformizer.
    \item Ramification index of a local extension.
    \item Residue degree of a local extension.
    \item Hensel's lemma.
\end{itemize}

\subsubsection*{Learning objectives}

Students should be able to:

\begin{itemize}
    \item Define and compare archimedean and nonarchimedean absolute values.
    \item Construct \(\mathbb{Q}_p\) as a completion of \(\mathbb{Q}\).
    \item Describe the topology induced by a nonarchimedean valuation.
    \item Use Hensel's lemma to lift factorizations or roots.
    \item Describe the ring of integers and residue field of a local field.
    \item Analyze finite extensions of local fields.
    \item Relate local ramification to global ramification.
\end{itemize}

\subsection{Module 5: Global Fields and the Local--Global Perspective}

\textbf{Suggested duration:} Weeks 9--10

\subsubsection*{Main themes}

\begin{itemize}
    \item Global fields.
    \item Places of a number field.
    \item Archimedean and nonarchimedean places.
    \item Completions at places.
    \item Local solvability and global solvability.
    \item Approximation theorems.
    \item Norms and local norms.
\end{itemize}

\subsubsection*{Suggested chapters}

\begin{enumerate}
    \item Global fields.
    \item Places and completions.
    \item Approximation and weak approximation.
    \item Local--global principles.
\end{enumerate}

\subsubsection*{Important concepts}

\begin{itemize}
    \item Global field.
    \item Place \(v\) of a number field.
    \item Completion \(K_v\).
    \item Product formula.
    \item Weak approximation.
    \item Strong approximation.
    \item Local solvability.
    \item Global solvability.
    \item Hasse-type principles.
\end{itemize}

For a number field \(K\), the product formula takes the form
\[
    \prod_v |x|_v^{n_v}=1,
    \qquad x\in K^\times,
\]
where the product is over all places \(v\) of \(K\), and the exponents
\(n_v\) depend on the normalization of the absolute values.

\subsubsection*{Learning objectives}

Students should be able to:

\begin{itemize}
    \item Identify the places of a number field.
    \item Explain the relation between a global field and its completions.
    \item State and apply approximation results.
    \item Formulate local--global questions precisely.
    \item Explain why local information can sometimes determine global
          arithmetic behavior.
\end{itemize}

\subsection{Module 6: Adèles and Idèles}

\textbf{Suggested duration:} Weeks 10--12

\subsubsection*{Main themes}

\begin{itemize}
    \item Restricted direct products.
    \item The ring of adèles.
    \item The group of idèles.
    \item Diagonal embeddings.
    \item Topological and algebraic structure.
    \item Class groups and idèle class groups.
    \item Global arithmetic through local components.
\end{itemize}

\subsubsection*{Suggested chapters}

\begin{enumerate}
    \item Restricted products of local fields.
    \item The ring of adèles.
    \item The idèle group.
    \item Idèle class groups.
    \item Applications to global arithmetic.
\end{enumerate}

\subsubsection*{Important concepts}

The ring of adèles of a number field \(K\) is the restricted product
\[
    \mathbb{A}_K
    =
    \prod_v' K_v,
\]
where, at all but finitely many nonarchimedean places \(v\), the component
belongs to the ring of integers \(\mathcal{O}_v\).

The group of idèles is
\[
    \mathbb{A}_K^\times
    =
    \prod_v' K_v^\times,
\]
with respect to the unit groups \(\mathcal{O}_v^\times\) at almost all
nonarchimedean places.

Important related objects include:

\begin{itemize}
    \item The diagonal embedding \(K\hookrightarrow\mathbb{A}_K\).
    \item The diagonal embedding \(K^\times\hookrightarrow\mathbb{A}_K^\times\).
    \item The idèle class group
          \(\mathbb{A}_K^\times/K^\times\).
    \item Restricted direct products.
    \item Local and global Haar measures.
    \item Characters of adèle groups.
\end{itemize}

\subsubsection*{Learning objectives}

Students should be able to:

\begin{itemize}
    \item Define restricted direct products rigorously.
    \item Construct the adèle ring of a number field.
    \item Explain the difference between adèles and idèles.
    \item Describe the diagonal embedding of a global field.
    \item Relate idèle class groups to ideal class groups.
    \item Explain how adèles unify archimedean and nonarchimedean data.
\end{itemize}

\subsection{Module 7: Tate's Thesis}

\textbf{Suggested duration:} Final module, if time permits

\subsubsection*{Main themes}

\begin{itemize}
    \item Harmonic analysis on local fields.
    \item Schwartz--Bruhat functions.
    \item Fourier transform over local fields.
    \item Characters of adèle groups.
    \item Poisson summation.
    \item Zeta integrals.
    \item Functional equations of \(L\)-functions.
\end{itemize}

\subsubsection*{Suggested chapters}

\begin{enumerate}
    \item Fourier analysis on local fields.
    \item Schwartz--Bruhat spaces.
    \item Poisson summation on adèles.
    \item Global zeta integrals.
    \item Functional equations and analytic continuation.
\end{enumerate}

\subsubsection*{Important concepts}

\begin{itemize}
    \item Schwartz space.
    \item Schwartz--Bruhat function.
    \item Additive and multiplicative characters.
    \item Fourier transform.
    \item Haar measure.
    \item Poisson summation formula.
    \item Tate zeta integral.
    \item Hecke \(L\)-function.
    \item Functional equation.
\end{itemize}

A typical global zeta integral has the schematic form
\[
    Z(\phi,\chi,s)
    =
    \int_{\mathbb{A}_K^\times}
    \phi(x)\chi(x)|x|^s\,d^\times x,
\]
where \(\phi\) is an appropriate Schwartz--Bruhat function and \(\chi\) is
a character of the idèle class group.

\subsubsection*{Learning objectives}

If this module is covered, students should be able to:

\begin{itemize}
    \item Explain the role of Fourier analysis in modern number theory.
    \item Describe Schwartz--Bruhat functions on local fields and adèles.
    \item State the adèlic Poisson summation formula.
    \item Define basic zeta integrals.
    \item Explain conceptually how functional equations arise from Fourier
          transform and Poisson summation.
\end{itemize}

\section{Proposed Weekly Breakdown}

\renewcommand{\arraystretch}{1.25}

\begin{longtable}{p{0.12\textwidth} p{0.28\textwidth} p{0.48\textwidth}}
\toprule
\textbf{Weeks} & \textbf{Unit} & \textbf{Representative topics} \\
\midrule
1--2 &
Algebraic preliminaries &
Number fields, rings of integers, integral elements, norms, traces,
embeddings, fractional ideals, localization. \\
\midrule
3--4 &
Dedekind domains &
Noetherian and integrally closed domains, ideal factorization, fractional
ideals, principal ideals, ideal class groups. \\
\midrule
5--7 &
Ramification &
Prime decomposition, residue degrees, ramification indices, discriminants,
differents, decomposition groups, inertia groups. \\
\midrule
7--9 &
Local fields &
Valuations, completions, \(\mathbb{Q}_p\), local rings, residue fields,
uniformizers, finite extensions, Hensel's lemma. \\
\midrule
9--10 &
Global and local fields &
Places, completions, product formula, approximation theorems, local--global
principles. \\
\midrule
10--12 &
Adèles and idèles &
Restricted products, adèle rings, idèle groups, diagonal embeddings,
idèle class groups, global arithmetic. \\
\midrule
12--14 &
Tate's thesis &
Fourier analysis, Haar measure, Schwartz--Bruhat functions, Poisson
summation, zeta integrals, functional equations, if time permits. \\
\bottomrule
\end{longtable}

\section{Major Difficult Topics}

The following topics are likely to require the greatest amount of abstract
reasoning and practice:

\subsection{Ideal Factorization versus Element Factorization}

Unique factorization may fail for elements in \(\mathcal{O}_K\), while
nonzero ideals in a Dedekind domain factor uniquely. Students must understand
why ideal factorization restores a useful form of arithmetic structure.

\subsection{Ramification}

Ramification combines ideal theory, residue fields, field extensions, and
Galois theory. The notation
\[
    \mathfrak{p}\mathcal{O}_L
    =
    \prod_{i=1}^{g}\mathfrak{P}_i^{e_i}
\]
must be interpreted simultaneously in terms of factorization, local
behavior, and extension degree.

\subsection{Discriminants and Differents}

Discriminants are defined using traces and determinants, but they also detect
ramification. Understanding the relationship between the discriminant ideal,
the different ideal, and ramified primes is conceptually demanding.

\subsection{Valuations and Completions}

Nonarchimedean geometry differs substantially from ordinary real analysis.
The ultrametric inequality, totally disconnected topology, and behavior of
convergent sequences require a change in intuition.

\subsection{Local--Global Principles}

A global arithmetic problem may be studied through all of its local
completions. However, local solvability does not automatically imply global
solvability in every setting. Students must distinguish valid local--global
theorems from situations in which obstructions occur.

\subsection{Restricted Products}

The definitions of adèles and idèles are compact but technically subtle.
Students must understand both the algebraic restricted product and the
topology it carries.

\subsection{Tate's Thesis}

Tate's thesis combines algebraic number theory, harmonic analysis, measure
theory, Fourier transform, and complex analysis. The main challenge is
understanding how local calculations assemble into a global functional
equation.

\section{Prerequisite Review Topics}

Students who need preparation should review the following material before or
during the first weeks of the course.

\subsection{Algebra}

\begin{itemize}
    \item Modules over principal ideal domains.
    \item Noetherian rings and modules.
    \item Localization.
    \item Integral extensions.
    \item Fraction fields.
    \item Tensor products.
    \item Finite and separable field extensions.
    \item Galois groups and the fundamental theorem of Galois theory.
\end{itemize}

\subsection{Commutative Algebra}

\begin{itemize}
    \item Prime and maximal ideals.
    \item Quotient rings.
    \item Factorization of ideals.
    \item Discrete valuation rings.
    \item Dimension theory.
    \item Integral closure.
    \item Dedekind domains.
\end{itemize}

\subsection{Analysis}

\begin{itemize}
    \item Metric and normed spaces.
    \item Completeness and compactness.
    \item Series and sequences.
    \item Measure and integration.
    \item Fourier series and Fourier transform.
    \item Basic complex analysis.
\end{itemize}

\subsection{Useful Mathematical Techniques}

\begin{itemize}
    \item Proof by contradiction and induction.
    \item Diagram chasing and homomorphism arguments.
    \item Computation with bases and matrices.
    \item Working with quotient structures.
    \item Translating between algebraic and topological definitions.
    \item Constructing examples and counterexamples.
\end{itemize}

\section{Homework and Assessment}

\subsection{Grading Policy}

The course grade is based entirely on homework:

\[
    \boxed{\text{Homework: }100\%}
\]

The lowest homework grade is dropped.

\subsection{Homework Format}

In general:

\begin{itemize}
    \item A problem set is posted on Canvas by the end of each Thursday.
    \item Solutions are due at the beginning of class the following Thursday.
    \item Late homework is not accepted without a valid excuse.
    \item One or two problems may be graded for accuracy.
    \item The remaining problems are generally graded for completion.
\end{itemize}

\subsection{Collaboration Policy}

Collaboration is allowed and encouraged after each student has made an
independent attempt. Every student must write the final solution individually
and in their own words. Direct copying from another student or another source
is prohibited.

\subsection{Recommended Homework Categories}

Homework should include a mixture of:

\begin{enumerate}
    \item Definitions and short proofs.
    \item Computations with rings of integers and ideals.
    \item Prime decomposition problems.
    \item Ramification and discriminant calculations.
    \item Valuation and completion exercises.
    \item Proofs involving local fields.
    \item Conceptual questions about adèles and idèles.
    \item Optional challenge problems.
\end{enumerate}

\section{Suggested Problem-Solving Progression}

Students should approach problems using the following progression:

\begin{enumerate}
    \item Identify the algebraic objects involved.
    \item Write down the relevant definitions precisely.
    \item Determine whether the problem is global, local, or both.
    \item Translate element questions into ideal-theoretic language when
          appropriate.
    \item Use norms, traces, valuations, or residue fields to obtain
          numerical information.
    \item Check whether a theorem applies and verify its hypotheses.
    \item Test the argument on a small example.
    \item Write a complete proof using clear notation and logical transitions.
\end{enumerate}

\section{Academic Integrity}

Students are expected to be honest and truthful in all course work. Prohibited
conduct includes:

\begin{itemize}
    \item Cheating.
    \item Plagiarism.
    \item Reuse of assignments.
    \item Unauthorized collaboration.
    \item Improper use of the Internet or electronic devices.
    \item Alteration of graded assignments.
    \item Forgery or falsification.
    \item Facilitating academic dishonesty.
\end{itemize}

Collaboration is permitted only under the conditions stated in the homework
policy. The submitted write-up must represent the student's own mathematical
communication.

\section{Accessibility and Student Support}

Students requiring accommodations should obtain an Accommodation Letter from
Student Disability Services and communicate with the instructor as early as
possible.

Student Disability Services may be contacted at:

\begin{itemize}
    \item Shaffer Hall \#101, Homewood Campus.
    \item Telephone: 410-516-4720.
    \item Email:
    \href{mailto:studentdisabilityservices@jhu.edu}
    {studentdisabilityservices@jhu.edu}.
\end{itemize}

Students experiencing anxiety, stress, depression, or other mental-health
concerns are encouraged to contact the JHU Counseling Center:

\begin{itemize}
    \item 3003 North Charles Street, Suite S-200.
    \item Telephone: 410-516-8278.
\end{itemize}

The course is intended to maintain a respectful and inclusive classroom
environment. Students should be treated with dignity regardless of identity,
background, or perspective.

\section{Summary of Course Architecture}

\begin{center}
\begin{tabular}{p{0.24\textwidth} p{0.29\textwidth} p{0.34\textwidth}}
\toprule
\textbf{Level} & \textbf{Primary focus} & \textbf{Expected outcome} \\
\midrule
Foundational &
Number fields, rings of integers, ideals &
Students can work with basic arithmetic objects in number fields. \\
Structural &
Dedekind domains and ideal classes &
Students understand unique ideal factorization and class groups. \\
Arithmetic &
Ramification and discriminants &
Students can analyze the behavior of primes in field extensions. \\
Local &
Valuations, completions, and local fields &
Students can transfer arithmetic questions to local settings. \\
Global &
Places, approximation, and local--global methods &
Students understand how local data interact with global arithmetic. \\
Advanced &
Adèles, idèles, and Tate's thesis &
Students see the modern unified language of algebraic number theory. \\
\bottomrule
\end{tabular}
\end{center}

\section{Overall Course Outcome}

The central goal of Math 617 is to develop a coherent understanding of
algebraic number theory from both global and local perspectives. By the end
of the course, students should be able to move between number fields, their
rings of integers, prime ideals, completions, and adèlic constructions.
They should also be able to read advanced mathematical texts and formulate
proofs involving the arithmetic structure of fields and rings.

\end{document}